The main experiments evaluate language-model reasoning on CLADDER, which generates hypothetical causal questions from structural causal models. The experiments used released benchmark questions and newly generated numerical instances. No participant recruitment, interviews, surveys or interventions on people formed part of the reported experiments. The results describe performance on constructed causal problems, which limits their use as evidence for decisions affecting real people.

A principal ethical concern is that a shorter explanation may retain a correct answer while omitting or misrepresenting a causal assumption. The evaluation therefore distinguishes answer accuracy, trace provenance and semantic validity. In the controlled few-shot evaluations, target prompts excluded their reference answers, and failed or malformed outputs remained in the reported denominators. Examples of answer-preserving reasoning errors are discussed because fluent explanations can otherwise invite unwarranted trust. The findings are not presented as validation of a system for consequential decision-making.

Qwen inference ran on rented AutoDL GPUs. Historical tool-agent experiments, shared-first-pass evaluation and semantic judging also used external GPT and DeepSeek services. These external computing and model services processed the research materials. Source questions, model responses and analysis records are stored separately for traceability. The author remains responsible for understanding and verifying the submitted work.
